\section{Solución propuesta}

Luego de lo mencionado anteriormente presentamos un pseudocodigo que muestra la solución propuesta por programación dinámica para este problema

\begin{figure}[H]
\centering
\begin{lstlisting}[style=fileformat]
Entrenamientos(e,s): 

	Sea N igual al largo de e o s mas 1
	Sea g una matriz de NxN con todos ceros
	Sea k un entero iniciado en 0
	Sea opt un arreglo de largo N inicializado con 0s
	
	Por cada i entre 1 a N:
		Por cada j entre 1 a N:
			Si i>j: 
				continuamos 
			Si j == i:
				g[i][j]=e[i]
			Sino:
				g[i][j]= g[i][j-1] + min(e[j],s[j-i])
	
	Por cada i entre 1 a N:
		Sea opt_k un arreglo de largo N inicializado en 0s
		Sea mayor_opt un entero inicializado en 0
			
		Por cada j entre 1 a i:
			Si j != i:
				opt_k[j] = opt[i-j-1]+g[i-j+1][i]
			Sino:
				opt_k[j] = g[1][j]
			
			Si opt_k[j]>=mayor_opt:
				mayor_opt = opt_k[j]
				k=j
		opt[i] = opt_k[k]

	devolver reconstruccion(opt, k)	
	
\end{lstlisting}
\label{cod:ordenado}
\caption{Pseudocodigó propuesto para resolver el problema}
\end{figure}


Siendo reconstrucción el algoritmo de para obtener la solución, que se vuelve mas simple guardando el ultimo k. Este algoritmo corresponde al siguiente pseudocodigo

\begin{figure}[H]
\centering
\begin{lstlisting}[style=fileformat]
Reconstruccion(opt,k): 

	Sea dia el calculo de len(opt)-1-k
	Por cada i entre 0 a k
		

    Sea rec una lista vacia con la solcuion
    Mientras len(rec) < len(opt):
         Sea dia igual a len(opt)-1 -k -1
         Por cada i entre 0 a k
            agregamos 'E' a rec
         agregamos 'D' a rec
         k=dia

    Damos vuelta la lista rec
    Borramos la primera componente que es la componente 'D' agregada al final
	Devolvemos rec

\end{lstlisting}
\label{cod:ordenado}
\caption{Pseudocodigó propuesto para reconstruir la solucion}
\end{figure}